\par\medskip
\hypertarget{native-086-article067}{}\pdfbookmark[2]{67.}{article-067}
\begin{center}67.\end{center}
Another method for solving the problem of the preceding article is derived from the equations\HistoricalNoteEight{Manuscript note:]
\[
\begin{aligned}
\sin\alpha\cos\delta&=-\sin\epsilon\sin b+\cos\epsilon\cos b\sin l\\[.35em]
\cos\alpha\cos\delta&=\cos b\cos l\\[.35em]
\sin\delta&=\cos\epsilon\sin b+\sin\epsilon\cos b\sin l.
\end{aligned}
\]
}
\[
\begin{aligned}
\cos\epsilon\sin l&=\sin\epsilon\operatorname{tang}b+\cos l\operatorname{tang}\alpha\\[.4em]
\sin\delta&=\cos\epsilon\sin b+\sin\epsilon\cos b\sin l\\[.4em]
\cos b\cos l&=\cos\alpha\cos\delta
\end{aligned}
\]
Let the auxiliary angle $\theta$ be determined by the equation
\[
\operatorname{tang}\theta=\frac{\operatorname{tang}b}{\sin l},
\]
and we have
\[
\begin{aligned}
\operatorname{tang}\alpha&=\frac{\cos(\epsilon+\theta)\operatorname{tang}l}{\cos\theta}\\[.6em]
\operatorname{tang}\delta&=\sin\alpha\operatorname{tang}(\epsilon+\theta),
\end{aligned}
\]
\SourceLeaf{087}To check the calculation, we may add to these equations
\[
\cos\delta=\frac{\cos b\cos l}{\cos\alpha}\quad\text{or}\quad
\cos\delta=\frac{\cos(\epsilon+\theta)\cos b\sin l}{\cos\theta\sin\alpha}\cdot
\]
The ambiguity in determining $\alpha$ by the second equation is removed by the fact that $\cos\alpha$ and $\cos l$ must have the same sign.

This method is less convenient if, in addition to $\alpha$ and $\delta$, we also need $E$. The most convenient formula for determining this angle is then $\cos E=\frac{\sin\epsilon\cos\alpha}{\cos b}=\frac{\sin\epsilon\cos l}{\cos\delta}\cdot$ But this formula cannot give $E$ accurately whenever $\pm\cos E$ differs little from unity. It also leaves it ambiguous whether $E$ is to be taken between $0$ and $180^0$ or between $180^0$ and $360^0$. The first drawback is seldom of any consequence, especially since the utmost precision in the value of $E$ is unnecessary for calculating the differential ratios. The ambiguity, however, is easily removed by the equation $\cos b\cos\delta\sin E=\cos\epsilon-\sin b\sin\delta$, which shows that $E$ must be taken between $0$ and $180^0$ or between $180^0$ and $360^0$ according as $\cos\epsilon$ is greater or less than $\sin b\sin\delta$. Clearly we need not even make this check whenever either of the angles $b$, $\delta$ does not exceed the limit $66^0\,32'$, for then $\sin E$ will always be positive. The same equation can also be used in the case indicated above to determine $E$ more accurately if it seems worth the trouble.
